\label{sec:methodology}

\subsection{Data Construction}

\paragraph{Incumbent identification.}
We identify all sitting members of the 117th Congress (2021--2022) who represent
districts in NC, WI, and TX, and who filed as candidates for reelection in the
2022 congressional elections following redistricting. Members who announced retirement
before redistricting was completed are excluded; members who sought different districts
due to pairing are included and classified as paired. This yields $n = 13$ incumbents
for NC (one member announced retirement before the filing deadline), $n = 8$ for WI
(all eight members filing), and $n = 36$ for TX (two members announced retirement
before filing).

\paragraph{Address geocoding.}
FEC Form 2 candidate filings for the 2022 election cycle report the candidate's home
address. We geocode these addresses to 2020 Census TIGER/Line census tracts using the
Census Bureau's online batch geocoder~\citep{censusGeocoder}. Where a home address
geocodes to a tract boundary (within 100 metres of a tract edge), we assign to the
tract containing the address centroid.

\paragraph{District assignment.}
For \textsc{bisect} plans, we run the algorithm with:
\begin{verbatim}
bisect build <state>_2022_test --year 2020
    --structure prime-factor
    --weights-override geographic
    --workers 8
\end{verbatim}
Single-seed results use seed\,$=42$ and are marked with $^\dagger$. The resulting
district assignment maps each census tract to a district identifier. We then look up
the district identifier for each incumbent's home tract to determine $d_P(t(i))$.

For enacted plans, we use the 2022 enacted congressional district shapefiles from
the Census Bureau's TIGER/Line redistricting layer, assigning each incumbent's
geocoded home address to the 2022 enacted district that contains it.

\subsection{Statistical Testing}

To test whether enacted pairing rates differ significantly from the random baseline,
we use a one-sided binomial test. Under the null hypothesis that each pair of
incumbents is paired with probability $p_0 = 1/k$ (or the geographically adjusted
$p_0 = \rho / k$), the observed pairing count $c$ follows a binomial distribution
with parameters $m = \binom{n}{2}$ (number of pairs) and $p_0$. The $p$-value for
the hypothesis that the enacted plan's pairing rate is below the random baseline is:
\[
  p\text{-value} = \Pr[C \leq c \mid C \sim \text{Binom}(m, p_0)].
\]

For \textsc{bisect} plans, we test the null hypothesis that the pairing rate is
consistent with the random baseline using the same binomial test, but two-sided:
\[
  p\text{-value} = 2 \cdot \min\left(\Pr[C \leq c], \Pr[C \geq c]\right).
\]
Failure to reject this null (i.e., $p > 0.10$) establishes that the \textsc{bisect}
pairing rate is consistent with random redistricting.

\subsection{Limitations}

Several limitations affect the interpretation of these results. First, the geocoding
of incumbent home addresses introduces measurement error: a small number of incumbents
may have moved since filing their FEC disclosures, and some addresses may geocode
to incorrect tracts. We estimate this affects at most 2--3 incumbents across the three
states and is unlikely to change any pairing outcome. Second, the random baseline
assumes that all geographically contiguous plans are equally likely, which is not
exactly true for \textsc{bisect}'s specific optimisation procedure. Third, incumbents
who lived in the same county or metropolitan area prior to redistricting may have been
more likely to be paired regardless of map, which the geographic adjustment partially
but not fully accounts for.
